\documentclass[whitelogo]{template/tudelft-report}
\usepackage[utf8]{inputenc}
\usepackage{changes}
\usepackage{graphicx}
\usepackage{adjustbox}
\usepackage{listings}

%bibliography imports


\begin{document}
%% Use Roman numerals for the page numbers of the title pages and table of
%% contents.
\frontmatter

%% Uncomment following 19 lines for a cover with a picture on the lower half only
%\title[tudelft-white]{Title}
%\subtitle[tudelft-cyan]{Optional subtitle}
%\author[tudelft-white]{J.\ Random Author}
%\affiliation{Technische Universiteit Delft}
%\coverimage{cover.jpg}
%\titleoffsetx{10cm}
%\titleoffsety{10cm}
%\afiloffsetx{1cm}
%\afiloffsety{18cm}
%\covertext[tudelft-white]{
%    \textbf{Cover Text} \\
%    possibly \\
%    spanning 
%    multiple 
%    lines
%    \vfill
%    ISBN 000-00-0000-000-0
%}
%\makecover

%% Uncomment following 16 lines for a cover with a picture on the lower half only
\title[tudelft-white]{Hall Effect}
\subtitle[tudelft-black]{Onderzoeken 5}
\author[tudelft-white]{Brian de Keijzer \\ Julia Norbart}
\affiliation{De Haagse Hogeschool}
\covertext[tudelft-white]{ % Dit moet makkelijker/mooier kunnen
    Vak \qquad\quad Onderzoeken 5 \\
    Klas\qquad\quad NH2B \\
    Docent \quad\space\space dr. ir. J. B. Oostinga \\
    \vfill
    \today
}
\setpagecolor{tudelft-cyan}
\makecover[split]



%% Include an optional title page.
\begin{titlepage}


\begin{center}

%% Insert the TU Delft logo at the bottom of the page.

%% Print the title in cyan.
{\makeatletter
\largetitlestyle\fontsize{64}{94}\selectfont\@title
%\largetitlestyle\color{tudelft-cyan}\Huge\@title
\makeatother}

%% Print the optional subtitle in black.
{\makeatletter
\ifx\@subtitle\undefined\else
    \bigskip
   {\tudsffamily\fontsize{22}{32}\selectfont\@subtitle}    
    %\titlefont\titleshape\LARGE\@subtitle
\fi
\makeatother}

\bigskip
\bigskip

by
%door

\bigskip
\bigskip

%% Print the name of the author.
{\makeatletter
%\largetitlefont\Large\bfseries\@author
\largetitlestyle\fontsize{26}{26}\selectfont\@author
\makeatother}

\bigskip
\bigskip

\vfill

\begin{tabular}{lll}
    Names: & Brian de Keijzer & Julia Norbart \\
    Studentnumbers: & 16011015 & 16083946 \\
    Emails: & 16011015@student.hhs.nl & 16083946@student.hhs.nl
\end{tabular}
%% Only include the following lines if confidentiality is applicable.

%\centering{\includegraphics{cover/logo_black}}


\end{center}

\begin{tikzpicture}[remember picture, overlay]
    \node at (current page.south)[anchor=south,inner sep=0pt]{
        \includegraphics{cover/logo_black}
    };
\end{tikzpicture}

\end{titlepage}



\chapter*{Summary}
\setheader{Summary}


\begin{flushright}
{\makeatletter\itshape
    \@author \\
    Utrecht, \today
\makeatother}
\end{flushright}

To get a better understanding of the forming scenarios and evolution of exoplanets, it is required to comprehend their atmospheric compositions. To do this, physically complex retrieval models are required. Retrievals with such complex models using traditional sampling methods usually takes hundreds to thousands of CPU core hours until convergence. The amount of computational power required creates a bottleneck for scientists. This will only be increasing with the amount of data coming available from future missions like ARIEL and JWST. The amount of time required for a retrieval can significantly be decreased by using machine learning methods like random forests and convolutional neural networks. A Generative Adversarial Network (GAN) is a type of neural network that has proven itself to be a reliable method for the analysis of complex physical phenomena. In recent work \cite{zingales2018exogan} a Deep Convolutional GAN (DCGAN) called ExoGAN has been developed, which together with semantic image inpainting performs retrievals in a matter of minutes. One of the drawbacks of ExoGAN is that the model has been trained on $10^9$ physically simple models containing no physics or chemistry. We propose SRON-GAN, a GAN which together with semantic image inpainting performs a retrieval on physically complex atmospheres in a matter of minutes. SRON-GAN is a DCGAN using a  Wasserstein GAN with Gradient Penalty (WGAN-GP) loss function. SRON-GAN will be trained on complex physical atmospheric models. These complex physical atmospheric models include molecular chemistry, self consistent temperature-pressure profile, and cloud formation. Once SRON-GAN is trained, it will be able to retrieve all the parameters it has been trained on from observed transmission spectra. All code is found on GitHub: https://github.com/deKeijzer/SRON-DCGAN .

\tableofcontents

%% Use Arabic numerals for the page numbers of the chapters.
\mainmatter

%% Import chapters
\chapter{Introduction}

The Hall voltage is a voltage due to the Hall effect. This effect is often use in modern equipment. In this report the Hall effect will be explained and used to measure the hall voltage. The Hall coefficient is defined as the ratio of the induced electric field to the product of the current density and the applied magnetic field. This coefficient depends on the material from which the conductor is made. In this experiment, the material that will be used to measure the Hall voltage is silver. 

\\
The research question is as follows:
 \begin{itemize}
     \item  What is the charge carrier density of silver?
 \end{itemize}
 
 




 

\chapter{Theory}
In this chapter the fundamental physics behind the Hall effect will be discussed. The Hall coefficient and charge carrier density will also be discussed. These both are unique material properties, which are related to each other. How so will become clear at the end of this chapter.

\section{The Hall effect}
When an electric current $I$ flows through a conductor perpendicular to a magnetic field $B$, it results in a transverse force $F_m$ on the moving charge carriers due to the magnetic field. Those charge carriers are pushed to one side of the conductor due to the Hall effect. All the chargers are building up at one side of the conductor. Because of this, there is a measurable voltage difference between the two sides of the conductor. This measurable voltage difference is called the Hall voltage. This Hall voltage is due is a potential difference due to the Hall effect. This effect is called after E.H. Hall, who discovered the Hall effect in 1897 [bron?]. The change in the direction of the current $I$ is caused by the magnetic force. This force is given by [bron?]:
    \begin{equation}
        \vec F = q\cdot (\vec E + \vec v \times \vec B)
        \label{eq:force}
    \end{equation}
%   
    \begin{table}[!htbp]
    \centering
    \caption{The quantities belonging to Equation \ref{eq:force}.}
    \label{table:force}
    \begin{tabular}{|c|c|c|}
    \hline
    \textbf{Quantity}   & \textbf{Description}           & \textbf{Unit}            \\ \hline
    $q$                  & Electrical charge    & {[}C{]}                     \\ \hline
    $\vec E$                  & Electric field strength         & {[}V/m{]}                     \\ \hline
    $\vec v$                  & Drift speed           & {[}m/s{]}                     \\ \hline
    $\vec B$                  & Magnetic field & {[}T{]} \\ \hline
    \end{tabular}
    \end{table}
    \begin{figure}[!htbp]
    \begin{center}
        \begin{minipage}[t]{0.45\textwidth}
            \includegraphics[scale=0.45]{figuren/theorie_e_rechtdoor.png}
            \caption{A schematic of the coil configuration.} \label{fig:schematic_setup_magnets}
        \end{minipage}
        \begin{minipage}[t]{0.45\textwidth}
            \includegraphics[scale=0.47]{figuren/theorie_e_afbuigen.png}
             \caption{The schematic setup for the experimental setup of the Hall effect.}\label{fig:schematic_hall_effect}
        \end{minipage}
    \end{center}
    \end{figure}
\begin{figure} [!htbp]
    \centering
    \includegraphics[scale=0.7]{figuren/schematic_forces.png}
    \caption{A schematic representation of the force on a current due to the Hall effect.}
    \label{fig:schematic_force}
\end{figure}
The quantities belonging to Equation \ref{eq:force} are defined in Table \ref{table:force}.
Because of this force, the positive and negative charges will have an acceleration in the direction of the force. This results in a change of the path of the charges. Which will lead to a potential difference.\\
Without magnetic field and thus without the force, the charges will go straight from one side of the material to the other, see Figure \ref{fig:schematic_setup_magnets}. When there is a magnetic field, the path changes in a curve, see Figure \ref{fig:schematic_hall_effect}. The potential different due to this effect is called the Hall Voltage. The Hall voltage is given by:
    \begin{equation}
        U_H = v_d\cdot B \cdot d
        \label{eq:Hall voltage}
    \end{equation}
With the distance $d$ between the positive and negative side of the plate. 
There also is current of the electrical charge in the negative direction with respect to the electon current, this is the conventional 'hole' current. This current can be written as:
    \begin{equation}
        I_x = n\cdot tw\cdot (-v_d)\cdot (-e)
        \label{eq:Hole current}
    \end{equation}
The quantities belonging to Equation \ref{eq:Hole current} are defined in Table \ref{table:Hole current}.
    \begin{table}[!htbp]
    \centering
    \caption{The quantities belonging to Equation \ref{eq:Hole current}.}
    \label{table:Hole current}
    \begin{tabular}{|c|c|c|}
    \hline
    \textbf{Quantity}   & \textbf{Description}           & \textbf{Unit}            \\ \hline
    $n$                  & Charge carrier density    & {[}Electrons/m^3{]}                     \\ \hline
    $tw$                  & The cross-sectional area         & {[}m^2{]}                     \\ \hline
    $-v_d$                  & Drift speed           & {[}m/s{]}                     \\ \hline
    $-e$                  & Charge of each electron & {[}C{]} \\ \hline
    \end{tabular}
    \end{table}
In which $e$ is the elementary charge, being equal to \cite{elementarycharge}:
\begin{equation}
    e = (1.602... \cdot 10^{-19} \pm 6.1 \cdot 10^{-9}) \quad \text{[C]} \label{eq:e}
\end{equation}
This equation can be solved for \emph{w} and plugged into Equation \ref{eq:Hall voltage} resulting in the Hall voltage:
 \begin{equation}
        U_H = \frac{1}{n \cdot e} \frac{B \cdot I}{d}. \label{eq:Hall Voltage2}
    \end{equation}

\section{Hall coefficient \& charge carrier density}
The Hall effect can be used to determine the Hall coefficient, which is an unique material property. this can be done by writing equation \ref{eq:Hall Voltage2} in the form of
    \begin{equation}
        U_H = R_H \frac{B I}{d} \label{eq:hallspanning_with_coefficient}
    \end{equation}
Where $R_H=\frac{1}{n \cdot e}$ is called the Hall coefficient. For silver $R_H$ is equal to \cite{apparatus_silver}:
    \begin{equation}
        \mid R_H \mid = (8,9 \pm 2) \cdot 10^{-11}  \quad \text{[m$^3$C$^{-1}$]} \label{eq:hallcoefficient_lit}
    \end{equation}
When the electric charge of the charge carriers in a material is known, Equation \ref{eq:hallspanning_with_coefficient} can be used to determine the charge carrier density. The charge carrier density $n$ is the number of charge carriers per volume.  The charge carrier density denotes the number of charge carriers per volume. It is measured in m$^{-3}$. As any density it can depend on position. It should not be confused with the charge density, which is the number of charges per volume at a given energy [bron?]. The charge carrier density is received by rewriting Equation \ref{eq:Hall Voltage2} to:
    \begin{equation}
        n = \frac{1}{e} \cdot \frac{1}{R_H}
    \end{equation}
In which $R_H$ is the Hall coefficient and $e$ the elementary charge. The charge carrier density cannot be directly measured. Therefor it has to be experimentally determined. This can be done by using the Hall effect. Silver for example has a charge carrier density of \cite{apparatus_silver}:
    \begin{equation}
        n = (6.6 \pm 2)\cdot 10^{28} \quad \text{[m$^{-3}$]} \label{eq:ccdensity_lit}
    \end{equation}
It is important to note that the uncertainties of Equation \ref{eq:hallcoefficient_lit} and Equation \ref{eq:ccdensity_lit} are not given by their source. However, their experimental results in comparison with the literature value are given. The difference between these two have been chosen as the uncertainties for Equation \ref{eq:hallcoefficient_lit} and Equation \ref{eq:ccdensity_lit}.
\chapter{Experimental procedure}
This chapter consists of the experimental setup and measurement procedure. Firstly the required apparatus is specified. After the apparatus is known, the calibration of the magnetic field is explained. At last the measurement procedure for the Hall voltage is discussed.

\section{Experimental setup}
The experimental setup can be divided into two sections. The first section is used to create a magnetic field. A schematic of this part of the experimental setup is shown in Figure \ref{fig:schematic_setup_magnets}. The maximum continues current supported by the coils is 5 A, this current is supplied by the power supply unit (PSU). It is possible to supply the coils with a higher current, for example 15 A. However, this has to be in a relatively short time frame. This is so the coils do not get damaged. A list of the required apparatus is shown in Figure \ref{fig:apparatus}, which can be found in Appendix-A. The power supply used to power the coils has to be able to deliver a current up to 5 A. This current is measured with a TTi 1604 digital multimeter. This multimeter is used for all current and voltage measurements. The necessary magnetic field strength to be created is between 0.1 and 0.6 T \cite{apparatus_silver}.\\
    \begin{figure}[!htbp]
    \begin{center}
        \begin{minipage}[t]{0.45\textwidth}
            \includegraphics[scale=0.35]{figuren/schematic_setup_magnets.png}
            \caption{A schematic of the coil configuration.} \label{fig:schematic_setup_magnets}
        \end{minipage}
        \begin{minipage}[t]{0.45\textwidth}
            \includegraphics[scale=0.5]{figuren/schematic_hall_voltage.png}
             \caption{The schematic setup for the experimental setup of the Hall effect.}\label{fig:schematic_hall_effect}
        \end{minipage}
    \end{center}
    \end{figure}
The second section of the experimental setup consists of the material on which the Hall voltage is measured. A schematic of this part of the experimental setup can be seen in Figure \ref{fig:schematic_hall_effect}. The PSU which supplies the current through the material must be able to supply up to 20 A. The Hall voltage is measured by a multimeter. An operational amplifier (Op-amp) might be needed to amplify the Hall voltage. The need for an Op-amp depends on the material used. Silver for example, has a Hall voltage of 5 $\mu$V when a current of 15 A flows through the material, in a 0,2 T magnetic field perpendicular to the flow of current \cite{halleffectzilver}. 5 $\mu$V cannot be measured by the TTi 1604 multimeter, therefor an Op-amp circuit is required to amplify the voltage to an order of magnitude which can be measured by the TTi 1604 \cite{dmm}. The chosen Op-amp is an AD521, which can amplify the signal up to one thousand times. With a maximum output voltage of 30 V \cite{AD521}. Note that the exact amplification of the Op-amp has to be determined before it is used in a measurement. This can be done by using two multimeters and a signal generator to apply a known signal and measure the Op-amps output. At last a computer with CASSY Lab is required for the data-acquisition of the Combi B-Sensor S. The Combi B-Sensor S measures the strength of the magnetic field at the Hall apparatus center with respect to earth's magnetic field.
    
\section{Magnetic field calibration}
The magnetic field strength created by the coils as a function of the current through them has to be known. This is due to the fact the magnetic field strength sensor does not fit in between the coils while the Hall apparatus is present. By experimenting it can be seen that the magnetic field in between the coils is homogeneous. For this reason a model for the magnetic field as function of current through the coils can respectively simply be made. This part of the experimental setup can be seen in Figure \ref{fig:cal_setup}.\\
    \begin{figure}[!htbp]
    \begin{center}
        \begin{minipage}[t]{0.45\textwidth}
            \includegraphics[scale=0.8]{figuren/schematic_calibration.png}
            \caption{Calibration of the magnetic field \\ schematically  \cite{halleffectzilver}.}\label{fig:cal_setup}
        \end{minipage}
        \begin{minipage}[t]{0.45\textwidth}
            \includegraphics[scale=0.8]{figuren/experimental_setup.png}
            \caption{The wiring diagram for the experimental setup of the Hall effect \cite{halleffectzilver}.}\label{fig:experimental_setup}
        \end{minipage}
    \end{center}
\end{figure}
The magnetic field as a function of the amount of current through the coils is calibrated as follows. The current through the coils is increased by steps of 0,1 A from 0 to 5 A. At each step the strength of the magnetic field is measured. This data can be used to create a mathematical model, which later on can be used to calculate the magnetic field strength when the current through the coils is known. This is done in Chapter 4.1. It is important to emphasize that the hysteresis characteristics of the soft metal surrounded by the coils is ignored due to simplicity.

\section{Hall voltage measurement}
The experimental setup should look like Figure \ref{fig:experimental_setup} when the coils and Hall apparatus are setup according to respectively Figure \ref{fig:schematic_setup_magnets} and Figure \ref{fig:schematic_hall_effect}. A Hall voltage can be create by two physical situations. When the current flow in the material is constant, the strength of the magnetic field perpendicular to this current flow has to increase. Or when the strength of the magnetic field is constant, the current has to increase. Either one will create a Hall voltage, as long as Equation \ref{eq:Hall Voltage2} is not equal to zero. Which means that $B$ and $I$ have to be larger then 0. \\
In case of the current through the coils being a constant, the magnetic field strength perpendicular to the material has to be increased. This is done by starting with 0 A through the coils, going up all the way to 5 A in steps of 0,5 A. This way the magnetic field $B$ can be calculated because the current through the coils is known, the Hall voltage at this specific magnetic field strength can also be measured. This should result in a linear model, since Equation \ref{eq:Hall Voltage2} is linear. \\
The same thing goes for a constant magnetic field strength of for example 0,2 T. The current through the material is varied from 0 to 20 A in steps of 1 A. This should also result in a linear model because of the reason stated above. \\
In this experiment a Hall apparatus with silver has been used in a 0,3 T magnetic field. It it important to note that two different silver apparatuses have been used with part number '58681 B2' \& '58681'. Why this is important will become clear in Chapter 6. When measuring the hall voltage, the offset voltage has to be determined first. This is done by turning off the magnetic field and measuring the offset voltage across the Hall apparatus. Once this is done the magnetic field can be turned on, which will result in a voltage increase due to the Hall effect. The actual difference between the measured voltage with magnetic field and the offset voltage is the voltage due to the Hall effect. One can also reverse the magnetic field to remove any longitudinal impurities in the force vector on the charge carriers.

\section{Python simulation}
Equation \ref{eq:Hall Voltage2} can be used to simulate experimental results. These simulations can be used to verify the actual experimental results while measuring.
    \begin{equation}
        U_H = \frac{1}{n \cdot e} \frac{B \cdot I}{d} \tag{\ref{eq:Hall Voltage2}}
    \end{equation}
Multiple physical situations have been modelled. In each situation all variables except for one, have been kept constant. These situations give insight in how the Hall voltage changes for different materials in different situations. In each situation Equation \ref{eq:e} has been used as the value for the elementary charge $e$.
\chapter{Measurements (Voorkeur naar eigen informatieve titel)}

Hoofdstuk 4 Resultaten: in dit hoofdstuk wordt de presentatie van de data gegeven in de vorm van grafieken en/of korte functionele tabellen. Leg uit wat je gemeten hebt, waar de data staat (in dit hoofdstuk / paragraaf of bijlage). Leg de grafieken uit en geef vooral aan wat de lezer uit de grafiek kan afleiden. Maak ook een voorbeeldberekening (inclusief de gebruikte formules) n.a.v. de data met direct daarna de onzekerheidsberekening. Bij een eenvoudige berekening (vb. wet van Ohm) hoeft geen voorbeeld berekening opgenomen te worden. De rest kan in een tabel of in een extra kolom bij de datatabel. Zijn er voor de grafieken al berekeningen nodig dan komen de berekeningen direct na de data presentatie. Bespreek de trendlijnen, pas regressieanalyse toe en breng de koppeling met de theorie aan. Schijf eerst tekst voordat je een tabel of grafiek presenteert.

Bij meetresultaten alleen relevante tabellen en/of grafieken (zie ook bijlagen) zodat de leesbaarheid niet verstoord wordt. Bij voorkeur een grafiek in plaats van een tabel. Een grafiek laat in één oogopslag het verloop van de relatie tussen de grootheden zien. De tabel moet wel in het verslag opgenomen worden; dit kan in de tekst (kleine tabel) of in een bijlage.

\section{Uncertainty Calculations}

\section{Data Analytics}
\chapter{Conclusion}
Hoofdstuk 5 Conclusie / discussie: de conclusie moet aansluiten op de inleiding. De lezer die de inleiding en de conclusie leest moet een helder beeld van het onderzoek /experiment gekregen hebben. In de conclusie worden antwoorden op de onderzoeksvragen en de resultaten van het experiment vergeleken met de literatuurwaarden of met je eigen verwachting. Werk in een conclusie altijd met getallen. Woorden als redelijk, matig, vrij goed, niet goed, enz. mogen niet erin voorkomen. Er zijn twee mogelijkheden: het resultaat van het experiment en de literatuurwaarde stemmen overeen of er is een (systematische) afwijking tussen het resultaat van het experiment en de literatuurwaarde. Zoek hiervoor een verklaring en geef commentaar. Aanbevelingen voor een vervolgonderzoek kunnen ook opgenomen worden. Geef geen nieuwe informatie in de conclusie en geen verwijzingen opnemen.
\chapter{Discussion}

Discussie: wanneer de discussie te lang voor de conclusie wordt dan kan na de conclusie een apart item Discussie opgenomen worden. In de discussie wordt de onderbouwing van je interpretatie van de resultaten opgenomen met daaruit voortvloeiend mogelijke aanbevelingen. Aanbevelingen kunnen zijn: verbetering meetopstelling, andere meetmethode, aanvullend onderzoek verrichten e.d.


\nocite{*}
\bibliographystyle{IEEEtran}
\bibliography{references}

%% Use letters for the chapter numbers of the appendices.
\appendix

\newpage

\chapter{Figures}

    \begin{figure}[!htbp]
    \begin{center}
    \includegraphics[scale=0.9]{figuren/apparatus.png}
    \end{center}
    \caption{A screenshot of the required apparatus \cite{halleffectzilver}.}\label{fig:apparatus}
    \end{figure}

    \begin{figure}[!htbp]
    \begin{center}
    \includegraphics[scale=0.9]{figuren/simulatie/B.png}
    \end{center}
    \caption{Simulation result of Formula \ref{eq:Hall Voltage2}.}  \label{fig:sim_B}
    \end{figure}
    
    \begin{figure}[!htbp]
    \begin{center}
    \includegraphics[scale=0.9]{figuren/simulatie/d.png} 
    \end{center}
    \caption{Simulation result of Formula \ref{eq:Hall Voltage2}.} \label{fig:sim_d}
    \end{figure}
    
    \begin{figure}[!htbp]
    \begin{center}
    \includegraphics[scale=0.9]{figuren/simulatie/I.png}
    \end{center}
    \caption{Simulation result of Formula \ref{eq:Hall Voltage2}.}  \label{fig:sim_I}
    \end{figure}

    \begin{figure}[!htbp]
    \begin{center}
    \includegraphics[scale=0.8]{figuren/resultaten/sample1.png}
    \end{center}
    \caption{Results sample 1 of the Hall voltage measurement with the Hall apparatus (silver).}\label{fig:silver1}
    \end{figure}
    

    \begin{figure}[!htbp]
    \begin{center}
    \includegraphics[scale=0.8]{figuren/resultaten/sample2.png}
    \end{center}
    \caption{Results sample 2 of the Hall voltage measurement with the Hall apparatus (silver).}\label{fig:silver2}
    \end{figure}
    
    \begin{figure}[!htbp]
    \begin{center}
    \includegraphics[scale=0.8]{figuren/resultaten/sample3.png}
    \end{center}
    \caption{Results sample 3 of the Hall voltage measurement with the Hall apparatus (silver).}\label{fig:silver3}
    \end{figure}
    
    \begin{figure}[!htbp]
    \begin{center}
    \includegraphics[scale=0.8]{figuren/resultaten/sample5.png}
    \end{center}
    \caption{Results sample 4 of the Hall voltage measurement with the Hall apparatus (silver).}\label{fig:silver4}
    \end{figure}
    
    \begin{figure}[!htbp]
    \begin{center}
    \includegraphics[scale=0.8]{figuren/resultaten/sample5.png}
    \end{center}
    \caption{Results sample 5 of the Hall voltage measurement with the Hall apparatus (silver).}\label{fig:silver5}
    \end{figure}
\input{appendix/python/pythonopmaak.tex} %vereist voor python opmaak
\chapter{Python code}
\section{Dataverwerking Life Time Meting}
\label{Bijlage2DataverwerkingLifeTimeMeting}
\lstinputlisting{appendix/python/Life_time_verwerking.py}

\section{Dataverwerking Delta Time Meting}
\label{Bijlage2DataverwerkingLifeTimeMeting}
\lstinputlisting{appendix/python/Delta_time_verwerking.py}
\chapter{Required appendix}
    Gezien dit de bijlage is en het twee onderdelen bevat welke als vereiste moeten worden toegevoegd aan de bijlage word dit onderdeel in het Nederlands geschreven.
    
\section{Verantwoording}
In deze sectie word omschreven wat de individuele bijdrages zijn geweest van beide duo-leden.
    \subsection{Brian de Keijzer}
    ...
    \subsection{Julia Norbart}
    ...

\section{Zelfreflectie}
Deze sectie bevat de zelfreflectie van beide duo-leden.
    \subsection{Brian de Keijzer}
    ...
    \subsection{Julia Norbart}
    ...



\end{document}